\chapter{Judd-Ofelt intensity parameters (menu 34)}
\label{ch:menu34}
\menulabel{34}

\section{Purpose}

The standard Judd-Ofelt (JO) theory describes the intensities of f-f
transitions as

\begin{equation}
  \mathcal{S}_{\mathrm{ED}} \; = \; \sum_{\lambda = 2, 4, 6}
    \Omega_\lambda \, \left| \langle \Psi_1 \| U^{(\lambda)} \| \Psi_2
    \rangle \right|^2 ,
\end{equation}

with the doubly-reduced unit-tensor matrix elements $U^{(\lambda)}$ being
host-independent table values of the ion, and the three parameters
$\Omega_\lambda$ obtained by least squares from measured line strengths.
This menu performs that fit and, from the fitted parameters, the emission
side (radiative rates, branching ratios, lifetime).  It is the landing
point of the transition-intensity data survey: the ORCA-side
intensity data (per-transition oscillator strengths and dipole moments of
CASSCF runs, with or without spin-orbit coupling) are read by menu 31 from
the property file's absorption section; the JO fit itself consumes the
literature conventions and tables, which this menu pins explicitly.

\section{What you need}

A dataset file (YAML; the schema is in the formats chapter): the host
refractive index (a constant, or Sellmeier terms), and per transition the
label, energy (or measured wavelength), oscillator strength, the lower
level's $J$, and the three $U^{(\lambda)}$ values.  Optional: an
$f_{\mathrm{md}}$ column (subtracted before the fit) and an emission block
(each entry additionally gives \code{j\_up}).

\section{Walkthrough}

The example fits the classic Eu$^{3+}$ benchmark (Babu et al., SET B --
the dataset of the fixture, with the $U^{(\lambda)}$ table of the 2024
literature) twice: unweighted, and with the normalized $1/S$ variant that
some authors use.  The two give different parameters from the same data --
which is exactly why the weighting is an explicit question.

\lstinputlisting[style=fbkcode, firstline=39, lastline=64,
  title={The menu-34 example (unweighted fit of the Eu$^{3+}$ benchmark:
  the parameters, the fit quality and the per-transition ratios)}]
  {examples/expected/m34_juddofelt.txt}

\section{What you get}

The fitted $\Omega_2/\Omega_4/\Omega_6$ (atomic units and the conventional
$10^{-20}$~cm$^2$), $\sigma$ and the relative $\sigma/\mathcal{S}_{\max}$,
the per-transition table ($S_{\exp}$, and $r = S_{\mathrm{ED}}/S_{\exp}$
for every fitted line), and -- when an emission block is present -- the
$A_{\mathrm{ED}}$ rates, branching ratios and the radiative lifetime.  The
report file \file{<dataset>.jo.fbk.md} carries the same content with the
full citations.

\section{Conventions and boundaries (all measured or source-verified)}

\begin{readbox}
\begin{itemize}
  \item \textbf{The local field is the squared virtual-cavity form}
    $\chi_{\mathrm{ED}} = (n_r^2+2)^2/9$.  The 2022/2024 source papers
    \emph{print} $(n_r^2+2)/9$ -- a missing square; the square is what
    their reference implementation uses
    (\code{calc\_chi\_ed(n,1) = (n**2+2)**2/9}) and what reproduces their
    published numbers.  This menu's regression test pins it: the Eu$^{3+}$
    benchmark reproduces the published relative $\sigma$ (8.52~\%) and
    $\Omega_6$ (2.253$\times10^{-20}$~cm$^2$) digit for digit.
  \item \textbf{Host dispersion}: a constant index, or Sellmeier
    $n_r^2(\lambda) = n_0^2 + \sum_i A_i\lambda^2/(\lambda^2 - B_i)$ with
    $B_i$ in $\mu$m$^2$ (wavelengths cut at 6000~nm, as the reference
    implementation); the measured wavelength drives both the Sellmeier
    lookup and the energy when no energy column is given.
  \item \textbf{A zero $U^{(\lambda)}$ column means that parameter is
    not determined} by the dataset; it is reported as such, never as a
    silent zero.
  \item \textbf{Magnetic-dipole parts are not computed} (they need
    free-ion wave functions): an optional \code{f\_md} column is
    subtracted before the fit, as the reference implementations do.
  \item \textbf{The extended $X_k$ / perturbative JO models} of the
    2022/2024 papers require an atomic-structure kernel (Cowan-style
    free-ion eigenvectors) and stay outside this toolkit's scope; the
    standard JO fit implemented here is what the sources themselves call
    ``the most recent experimental studies still use''.
\end{itemize}
\end{readbox}

\section{The ORCA side of transition intensities}

For \emph{ab initio} intensities, an ORCA state-averaged CASSCF job (with or
without \code{rel DoSOC}) prints the one-photon absorption spectrum --
per transition: energy, wavelength, oscillator strength \code{fosc(D2)},
dipole strength $D^2$, and the dipole components; the property file's
\code{\$CASSCF\_Absorption\_Spectrum} carries the same numbers in eleven
typed columns, measured as
$[\,\mathrm{eV}, \mathrm{cm}^{-1}, \mathrm{nm}, f_{\mathrm{osc}}, D^2,
D_X^{(\mathrm{re,im})}, D_Y^{(\mathrm{re,im})}, D_Z^{(\mathrm{re,im})}\,]$,
with two internal gates (the oscillator strength against
$\tfrac23\,\Delta E_{\mathrm{a.u.}} D^2$, and $D^2$ against the squared
components) that reproduce the printed numbers to their precision.  A job
with spin-orbit coupling prints the section twice; the copy with the
higher \code{\&RelCorrection} (the SOC-corrected one) is the one read.  The
magnetic-dipole data live in the separate \code{\$CASSCF\_ECD\_Spectrum}
section.  Mapping these per-level intensities onto the $J$-multiplet JO
parameters needs the $U^{(\lambda)}$ table (or crystal-field-level JO
theory, a registered candidate increment) -- which is why this menu's input
is a dataset file, not an ORCA output.
